\section{Nutritional self-sufficiency analysis}

\subsection{The calorie problem}

Full caloric independence from a small greenhouse is not achievable for most people with the crops CEA does best. The caloric yield per unit area of the typical CEA crop profile is too low \cite{patternbase2024calories}:

\begin{center}
\begin{tabular}{lcc}
\toprule
\textbf{Crop} & \textbf{Yield (kcal/m\textsuperscript{2}/year)} & \textbf{Area for one adult (m\textsuperscript{2})} \\
\midrule
Lettuce       & 300   & 2,400 \\
Spinach       & 500   & 1,460 \\
Tomato        & 3,000 & 243   \\
Dry beans     & 4,000 & 183   \\
Potato        & 7,000 & 104   \\
Sweet potato  & 8,000 & 91    \\
\bottomrule
\end{tabular}
\end{center}

An adult requires approximately 730,000~kcal/year. A mixed CEA system optimised for calorie-dense crops (sweet potato, potato, dry beans) could cover one adult's caloric needs in roughly 100--200~m\textsuperscript{2} of greenhouse space --- feasible for a large backyard installation but not a typical small greenhouse.

However, the crop table above understates the caloric potential of CEA by omitting two categories: mushrooms and oilseed legumes. Mushrooms are already a mainstream CEA crop --- most commercial production is controlled-environment by definition --- and require no light, reducing energy costs substantially. Dried mushrooms are storable and protein-dense. Oilseed legumes such as peanuts (~5,670~kcal/kg) and soybeans (~4,460~kcal/kg) exceed even wheat flour in caloric density and can be processed into oils, flours, and pastes with long shelf lives. Similarly, potatoes and sweet potatoes can be processed into storable flours and starches, extending their utility beyond fresh consumption.

The feasibility of these crops in CEA varies. Mushrooms and tubers are well-established; dry beans and soybeans have existing controlled-environment research; peanuts and other oilseed crops (sunflower, sesame) are plausible but largely unproven in CEA and represent R\&D targets for the project. The species recommendation engine (Section~[TODO]) should classify crops by evidence tier --- demonstrated, pilot-scale, and experimental --- and prioritise accordingly while the R\&D programme expands the proven portfolio.

\subsection{The correct framing: nutritional self-sufficiency}

The appropriate goal is not caloric independence but \textbf{nutritional self-sufficiency} --- specifically, the ability to produce the full spectrum of macronutrients (carbohydrates, proteins, fats) and micronutrients (vitamins, minerals, phytonutrients) needed for health, using a diversified crop portfolio that includes calorie-dense tubers, protein- and fat-contributing legumes and mushrooms, and micronutrient-dense vegetables and herbs.

This framing is both more honest and more compelling. It is the micronutrient and phytonutrient density of fresh produce that the industrial supply chain does worst at preserving --- nutrients degrade during the long transport and storage times that are inherent to centralised production. A locally grown, freshly harvested crop is nutritionally superior to the same crop shipped across a continent.

This objective also connects directly to the dietary shift documented in Section~\ref{sec:nutrition}: the industrial food system's narrowing of crop diversity toward a few commodity grains and oilseeds has driven the dominance of ultra-processed foods and the decline of micronutrient density in staple crops. CEA does not replace commodity grain production, but it can supply the diverse fresh produce, legumes, and tubers that a healthier dietary pattern requires --- making the dietary shift locally achievable rather than dependent on the same centralised supply chains whose fragility the problem analysis documents.

The species recommendation engine should therefore optimise for \textbf{nutritional completeness} as a primary objective, with caloric contribution and energy cost as co-objectives. This makes the goal achievable in realistic backyard footprints for the majority of a person's fresh produce needs.

\subsection{Vertical stacking multiplier}

CEA systems can stack crops vertically, multiplying the effective growing area per unit of floor space by 10--20$\times$. This means the effective growing area available in even a modest installation can substantially exceed its floor footprint, making the nutritional self-sufficiency target considerably more achievable than the raw floor-area numbers suggest.
